%=============================================================================
% WAVE 4, ITERATION 17: P1+P8 DEAD PATENTS + CROSS-PATENT ERROR CHECK
%
% Agent: W4i17 Agent 1 (Opus 4.6)
% Date: 20 March 2026
%
% MANDATE:
%   1. Confirm P1 and P8 DEAD (1 paragraph each)
%   2. Cross-check ALL i16 findings for internal consistency
%   3. RESOLVE the c_psi CONTRADICTION between i16 Agent 4 and Agent 2
%   4. Assess whether the corpus rewrite correctly reflects i15 state
%   5. Determine: P12 revival status, SQMS bound reliability, detection programme
%
% Status flags: [VERIFIED], [NEW], [CORRECTED], [CAUTION], [HONEST]
%=============================================================================

\documentclass[11pt]{article}
\usepackage{amsmath,amssymb,amsthm}
\usepackage{geometry}
\geometry{margin=1in}
\usepackage{booktabs}
\usepackage{enumitem}
\usepackage{hyperref}
\usepackage{xcolor}
\usepackage{tcolorbox}
\usepackage{float}
\usepackage{siunitx}
\usepackage{array}

\newtheorem{theorem}{Theorem}
\newtheorem{finding}{Finding}
\newtheorem{correction}{Correction}
\newtheorem{verdict}{Verdict}

\newcommand{\ee}[1]{\times 10^{#1}}
\newcommand{\psifield}{\psi}
\newcommand{\dpsi}{\delta\psi}
\newcommand{\dpsiEM}{\delta\psi_{\rm EM}}
\newcommand{\lbone}{|\lambda - 1|}
\newcommand{\Qpsi}{Q_\psi}
\newcommand{\Opsi}{\Omega_\psi}

\title{\textbf{Wave 4 Iteration 17: P1+P8 Dead Patent Confirmation\\
and Cross-Patent Error Check}\\[8pt]
{\large Resolution of the $c_\psi$ Contradiction Between i16 Agents}\\[4pt]
{\normalsize TOP 5 FINDINGS --- PROGRAMME-CRITICAL}}
\author{DFD Research Programme --- Agent 1 (W4i17, Opus 4.6)}
\date{20 March 2026}

\begin{document}
\maketitle

%=============================================================================
\section*{P1 and P8: Death Confirmation}
%=============================================================================

\textbf{P1 --- Drag Reduction / Gravitational-Doppler Sensing: DEAD.}
P1 has been dead since W3i7 (confirmed through 6 consecutive iterations, W3i7--W4i16). Every pathway has been killed: drag reduction requires $\lbone > 1.85\ee{-2}$ (exceeding bounds by $\sim 4\ee{4}\times$); Kerr-analog frame-dragging is $10^{38}$ below GR; GPS detection requires $\lbone > 7\ee{8}$; the P1+P5 bubble timing signal (4.0~ns) was killed at W3i5 because the required EM energy density ($5.8\ee{16}$~J/m$^3$) exceeds SRF capability by $10^{12}\times$, reducing the real signal to $\sim 0.81$~zs. The $c^2/2$ lever arm that makes tiny $\psi$-gradients produce large forces also means generating those gradients requires inaccessible EM energy densities. Zero resources recommended. Assessment is FINAL.

\textbf{P8 --- Deep-Space $\psi$-Communications: DEAD.}
P8 has been dead since W3i7 (confirmed through all subsequent iterations). The minimum transmitter mass for 1~kbit/s at 1~AU is 1.6~Earth masses. All modulation mechanisms fail: $\psi$-field propagation at $c$ is confirmed and inherent secrecy is real, but the gravitational coupling ($G/c^4$) makes the conversion efficiency catastrophically low. DSN wins by $370\times$ on Mars rate. Submarine communications (the last candidate application) requires Ceres-mass transmitters. No resolution pathway exists. Zero resources recommended. Assessment is FINAL.

\newpage

%=============================================================================
\section*{Executive Summary: The $c_\psi$ Contradiction Resolution}
%=============================================================================

\begin{tcolorbox}[colback=red!5!white, colframe=red!60!black,
  title=\textbf{W4i17 TOP 5 FINDINGS}]

\begin{enumerate}

\item \textbf{FINDING 1 [PROGRAMME-CRITICAL]: The i16 agents are NOT
  contradicting each other --- they are evaluating different physical
  quantities. BOTH are correct in their respective domains.}
  Agent~4 (c$_\psi$ audit) correctly identifies that for large-$\Delta$
  perturbations ($\Delta \gg 1$), the verification equation (\S2.6)
  applies and $c_\psi = c$. Agent~2 ($Q_\psi$ analysis) correctly
  derives $\tilde{c}_\psi \approx 10^{-13}$~m/s from $K''(\Delta_{\rm bg})$
  at $\Delta_{\rm bg} \sim 2.5\ee{5}$. The discrepancy arises because
  Agent~2 evaluates the propagation speed of \emph{small perturbations
  about the laboratory background}, where $K''(\Delta_{\rm bg}) \approx
  1/\Delta_{\rm bg}^2 \approx 1.6\ee{-11}$, while Agent~4 evaluates
  the propagation speed of the perturbation \emph{itself} when $\Delta
  \gg 1$. See Finding~1 for the full resolution.

\item \textbf{FINDING 2 [PROGRAMME-CRITICAL]: The correct $c_\psi$ for
  SRF cavity psi-modes is $c$, not $10^{-13}$~m/s. Agent~4 is right
  for the detection programme.} The perturbation $\dpsiEM$ driven by
  the SRF EM field has its own $\Delta$ given by
  $\Delta_{\rm pert} = (c/a_0)\omega A_\psi \sim 10^5$ (at $|\lambda-1|
  \sim 10^{-9}$). Since $\Delta_{\rm pert} \gg 1$, the perturbation
  is in the large-$\Delta$ regime where $c_\psi = c$. Agent~2's
  $10^{-13}$~m/s applies to perturbations whose \emph{own} $\Delta$ is
  set by the background, not the perturbation's frequency content.
  P12 STAYS DEAD (but for reasons unrelated to $c_\psi$).
  SQMS bound is RELIABLE (with caveats).

\item \textbf{FINDING 3 [CORRECTED]: Agent~2 made a linearization error
  --- evaluating $K''$ at the BACKGROUND $\Delta_{\rm bg}$ rather than
  at the PERTURBATION's $\Delta$.} In the DFD temporal sector, the
  absolute value $|\dot\psi - \dot\psi_0|$ means the relevant $\Delta$
  for linearization is the perturbation's own temporal invariant, not
  the background's. For a perturbation oscillating at frequency $\omega$
  with amplitude $A_\psi$, $\Delta_{\rm pert} = (c/a_0)\omega A_\psi$.
  Agent~2 correctly computed $\Delta_{\rm bg} = 2.5\ee{5}$ from
  Earth's rotation but then used $K''(\Delta_{\rm bg})$ in the wave
  equation coefficient. The correct procedure: linearize about
  $\psi_{\rm grav}$ with $\dot\psi_0 \neq 0$, so the perturbation's
  $\delta\Delta = (c/a_0)|\delta\dot\psi|$. For $\delta\Delta \gg 1$
  (laboratory), $K''(\delta\Delta) \sim 1/(\delta\Delta)^2$, and the
  full wave equation recovers $c_\psi = c$.

\item \textbf{FINDING 4 [CAUTION]: The SQMS bound $\lbone < 1.56\ee{-9}$
  is RELIABLE but carries a $\sqrt{\Qpsi}$ uncertainty that Agent~2
  correctly flagged.} With $c_\psi = c$ confirmed for laboratory
  perturbations, the psi-mode frequency $\Opsi = c k_{\rm cav} \sim
  10^{10}$~rad/s and the loss formula from W4i15 applies. However,
  the \emph{local cavity} $\Qpsi$ (not the MOND $10^9$) should enter.
  The radiation-$Q$ estimate $\Qpsi^{\rm local} \sim (kR)^3 \sim 31$
  is now restored (since $c_\psi = c$ validates the radiation escape
  assumption). This would TIGHTEN the SQMS bound by $\sqrt{10^9/31}
  \approx 5700\times$ to $\lbone < 2.7\ee{-13}$. However, whether the
  cavity walls partially confine the psi-mode (enhancing $\Qpsi^{\rm local}$
  above 31) remains an open theoretical question. Conservative range:
  $\Qpsi^{\rm local} \in [31, 10^9]$, giving $\lbone < [2.7\ee{-13},
  1.56\ee{-9}]$. The Q-ratio shape diagnostic ($\mathcal{R} = 0.49$)
  is ROBUST ($\Qpsi$ cancels).

\item \textbf{FINDING 5 [CORRECTED]: The corpus rewrite contains 4
  errors relative to the i16 findings it should reflect.}
  (a)~P12 is listed as ``DEAD'' based on $c_\psi \sim 10^{-5}$~m/s,
  but i16 Agent~4 showed $c_\psi = c$ in the lab. P12's death should
  be reviewed --- the single-cavity $C_1 = 0.47$ survives, and
  multi-cavity coherence is restored if $c_\psi = c$. However,
  multi-cavity superradiance may still fail for OTHER reasons
  (inter-cavity coupling geometry, mode matching).
  (b)~The corpus states ``Dust-branch $w \sim 10^{-37}$'' (two
  locations), but i16 corrected this to $w \sim 5\ee{-6}$.
  (c)~The corpus states ``$c_\psi \sim 10^{-5}$~m/s (Dust branch,
  W4i14)'' in the key quantities table. This should be
  ``$c_\psi = c$ (laboratory, large-$\Delta$); $c_\psi \ll c$
  (cosmological, small-$\Delta$); regime-dependent (W4i16).''
  (d)~The ``superluminal $\psi$-signaling'' entry in the Definitively
  Closed list says ``$\psi$ propagates at $c_\psi \sim 10^{-5}$~m/s,
  NOT $c$'' --- this is wrong for laboratory perturbations where
  $c_\psi = c$.

\end{enumerate}
\end{tcolorbox}


%=============================================================================
%=============================================================================
\part{Finding 1: Resolution of the $c_\psi$ Contradiction}
\label{part:contradiction}
%=============================================================================
%=============================================================================

%=============================================================================
\section{The Two Agents' Derivations, Step by Step}
%=============================================================================

\subsection{Agent~4 (W4i16 c$_\psi$ Audit): $c_\psi = c$ in Lab}

Agent~4's argument proceeds as follows:

\begin{enumerate}
\item The full dynamic action (section02 Eq.~(6), reproduced below)
  contains spatial kinetic function $W(y)$ with $y = |\nabla\psi|^2/a_*^2$
  and temporal kinetic function $K(\Delta)$ with $\Delta =
  (c/a_0)|\dot\psi - \dot\psi_0|$:
  \begin{equation}
  S_\psi = \int dt\,d^3x\;\left\{
    \frac{a_*^2}{8\pi G}\left[
      W(y) + K(\Delta)
    \right] - \frac{c^2}{2}\psi(\rho - \bar\rho)
  \right\}.
  \end{equation}

\item Agent~4 identifies a normalization discrepancy: the action prefactor
  $a_*^2/(8\pi G)$ has dimensions kg$\cdot$s$^2$/m$^5$, while the matter
  coupling $c^2\psi\rho/2$ has dimensions kg/(m$\cdot$s$^2$). These
  differ by $c^4/4$. The field equation (Eq.~(8)) is correct; the action
  prefactor needs a $c^4/4$ correction.

\item Agent~4 then computes $\Delta$ for a laboratory perturbation:
  at $\omega \sim 6$~GHz, $A_\psi \sim |\lambda-1| \sim 10^{-24}$
  (or even $10^{-9}$ at the SQMS bound), $\Delta \sim (c/a_0)\omega
  A_\psi \sim 2.5\ee{18} \times 4\ee{10} \times 10^{-24} \sim 10^5$.

\item Since $\Delta \gg 1$, $K'(\Delta) \approx 1$, and the linearized
  equation reduces to the \S2.6 verification equation:
  $\nabla^2\psi - (1/c^2)\ddot\psi = -8\pi G\rho/c^2$.
  Propagation speed $= c$.

\item For cosmological perturbations ($\omega \sim H_0$, $A_\psi \sim
  10^{-5}$), $\Delta \sim 5\ee{-5} \ll 1$, and the dust branch applies:
  $c_s \to 0$. The propagation speed is regime-dependent, not universal.
\end{enumerate}

\textbf{Agent~4 conclusion}: $c_\psi = c$ in the laboratory. The i14
result $c_\psi \sim 10^{-5}$~m/s applies only to cosmological
perturbations.

\subsection{Agent~2 (W4i16 $Q_\psi$ Local vs.\ Global): $c_\psi \sim 10^{-13}$~m/s}

Agent~2's argument proceeds differently:

\begin{enumerate}
\item Agent~2 starts from the linearized perturbation equation for
  $\dpsiEM$ about the gravitational background $\psi_{\rm grav}$:
  \begin{equation}
  \frac{a_*^2}{8\pi G}\left[
    \nabla^2(\dpsiEM) - K''(\Delta_{\rm bg})\,\frac{c^2}{a_0^2}\,
    \ddot{\delta\psi}_{\rm EM}
  \right] = -\frac{\lambda\,4\pi G}{c^4}\,u_{\rm EM}.
  \label{eq:agent2-eom}
  \end{equation}

\item Agent~2 evaluates $\Delta_{\rm bg}$ at the laboratory:
  $\dot\psi_{\rm local} \sim (2g_\oplus/c^2)v_{\rm rot} \sim 10^{-13}$~s$^{-1}$,
  giving $\Delta_{\rm bg} = (c/a_0)\dot\psi_{\rm local} \sim 2.5\ee{5}$.

\item Then $K''(\Delta_{\rm bg}) = 1/(1+\Delta_{\rm bg})^2 \approx 1.6\ee{-11}$.

\item The effective propagation speed:
  $\tilde{c}_\psi^2 = a_0^2/(c^2 K''(\Delta_{\rm bg}))
  = a_0^2\Delta_{\rm bg}^2/c^2 \approx 10^{-26}$~m$^2$/s$^2$,
  $\tilde{c}_\psi \approx 10^{-13}$~m/s.

\item The SRF drive frequency ($2\omega \sim 10^{10}$~rad/s) is
  $10^{22}\times$ above $\Opsi = \tilde{c}_\psi k_{\rm cav} \sim
  3\ee{-12}$~rad/s, putting the cavity deep in the
  $\omega \gg \Opsi$ regime. The dissipation is suppressed by
  $(\Opsi/\omega)^3 \sim 10^{-66}$.
\end{enumerate}

\textbf{Agent~2 conclusion}: $c_\psi \sim 10^{-13}$~m/s in the
laboratory. The SQMS psi-channel loss is suppressed by $\sim 10^{-66}$
relative to prior estimates. The entire Q-ratio detection programme
may be invalidated.


%=============================================================================
\section{Where Agent~2 Goes Wrong}
\label{sec:error}
%=============================================================================

\begin{finding}[The linearization error]
\label{find:linearization}
Agent~2's error is in step~1 above: the coefficient $K''(\Delta_{\rm bg})$
in Eq.~\eqref{eq:agent2-eom}.

The temporal kinetic function $K(\Delta)$ with $\Delta = (c/a_0)|\dot\psi
- \dot\psi_0|$ is a function of the \textbf{total} field's temporal
deviation, not the background alone. When we write $\psi = \psi_{\rm grav}
+ \dpsiEM$ and linearize, the argument of $K$ is:
\begin{equation}
\Delta = \frac{c}{a_0}|\dot\psi_{\rm grav} + \dot{\delta\psi}_{\rm EM}
- \dot\psi_0|.
\end{equation}

The perturbation $\dot{\delta\psi}_{\rm EM}$ oscillates at the SRF
frequency $\omega$ with amplitude $\sim |\lambda-1| \omega$. Agent~2
correctly computes $\Delta_{\rm bg} = (c/a_0)|\dot\psi_{\rm grav} -
\dot\psi_0| \sim 2.5\ee{5}$. But the perturbation's own contribution
to $\Delta$ is:
\begin{equation}
\delta\Delta = \frac{c}{a_0}|\dot{\delta\psi}_{\rm EM}|
\sim \frac{c}{a_0}\,\omega\,A_\psi
\sim 2.5\ee{18} \times 4\ee{10} \times A_\psi.
\end{equation}

For $A_\psi = |\lambda - 1| \sim 10^{-9}$ (SQMS bound):
$\delta\Delta \sim 10^{20}$. For $A_\psi \sim 10^{-24}$:
$\delta\Delta \sim 10^5$.

In either case, $\delta\Delta \gg \Delta_{\rm bg} \sim 10^5$ (at the
SQMS bound) or $\delta\Delta \sim \Delta_{\rm bg}$ (at extreme
suppression). The key point: \textbf{the perturbation is not ``small''
in the $\Delta$ sense}. The temporal kinetic function responds to the
full $\Delta$ including the perturbation itself.

This means standard linearization (expanding $K$ to quadratic order
about $\Delta_{\rm bg}$) is not valid when $\delta\Delta \gtrsim
\Delta_{\rm bg}$, which is the case for all laboratory-frequency
perturbations. Agent~4 recognized this (their ``key realization'':
the linearized equation is amplitude-dependent because of the absolute
value in $\Delta$).

The correct procedure for a perturbation with $\delta\Delta \gg 1$
is to note that $K'(\Delta) \to 1$ for $\Delta \gg 1$, so the temporal
sector acts as a standard $(1/c^2)\ddot\psi$ term, giving $c_\psi = c$.
This is precisely what Agent~4 derives.

Agent~2's error: they evaluated $K''$ at $\Delta_{\rm bg}$ (the
background's temporal invariant) rather than at the total
$\Delta_{\rm bg} + \delta\Delta$ (which is dominated by
$\delta\Delta$ at SRF frequencies). This produced an artificially
small $K''$ value ($1.6\ee{-11}$), which inflated the temporal
coefficient and suppressed $c_\psi$.

\textbf{Verdict}: Agent~4 is correct. $c_\psi = c$ for all
laboratory-frequency, SRF-amplitude perturbations. Agent~2's
$c_\psi \sim 10^{-13}$~m/s would apply only to perturbations
with $\delta\Delta \ll 1$, i.e., perturbations with
$\omega A_\psi \ll a_0/c \sim 4\ee{-19}$~s$^{-1}$. At SRF
frequencies ($\omega \sim 10^{10}$), this requires $A_\psi \ll
4\ee{-29}$ --- far below any physically relevant amplitude.
\end{finding}


%=============================================================================
\section{Why Agent~2's Honest Assessment Was Correct in Spirit}
\label{sec:honest}
%=============================================================================

Agent~2 flagged (Finding~5, ``The Deep Question'') that the temporal
sector is the key uncertainty and that the dust-branch theorem was
derived for FRW backgrounds. Their scenarios A/B/C analysis correctly
identified that the resolution depends on which $\Delta$ regime
applies. The error was in evaluating Scenario~C (laboratory $\Delta$)
using $\Delta_{\rm bg}$ rather than $\Delta_{\rm bg} + \delta\Delta$.

Agent~2 also correctly noted that the Q-ratio shape diagnostic
($\mathcal{R} = 0.49$) is robust against all scenarios, because
$\Qpsi$ cancels in the ratio. This finding stands.


%=============================================================================
%=============================================================================
\part{Finding 2: Downstream Consequences}
\label{part:downstream}
%=============================================================================
%=============================================================================

%=============================================================================
\section{P12 Status: DEAD Confirmed, But for Different Reasons}
%=============================================================================

With $c_\psi = c$ restored for laboratory perturbations, the i15
argument that killed P12 (coherence length $5\ee{-17}$~m at
$c_\psi \sim 10^{-5}$~m/s) is invalidated. The coherence length
is restored to $\ell_{\rm coh} = c \Qpsi/(2\pi f) \sim 1$~m at
$\Qpsi \sim 10^2$.

However, P12 multi-cavity superradiance faces other obstacles that
were not resolved by the $c_\psi$ correction:
\begin{enumerate}
\item \textbf{Single-pass conversion}: Degraded $10^{25}\times$ (W3i5).
  Single-cavity cooperativity $C_1 = 0.47$ requires $|\lambda-1| \sim 1$
  (at the current SRF parameters), which is excluded by the SQMS bound.
\item \textbf{Inter-cavity mode matching}: Multi-cavity Dicke
  superradiance was already corrected from $N^2$ to $N C_1$ (W4i11).
  Even with $N = 3$ cavities, $C_N = 3 \times C_1$ does not overcome
  the fundamental coupling weakness.
\item \textbf{The cooperativity at $|\lambda-1| \sim 10^{-9}$}:
  $C_1 \propto |\lambda-1|^2$, so at the SQMS bound,
  $C_1 \sim 10^{-18} \times C_1^{(\lambda=2)} \ll 1$.
\end{enumerate}

\textbf{Verdict}: P12 STAYS DEAD. The death cause should be updated
from ``$c_\psi$ kills coherence'' to ``cooperativity too low at
allowed $|\lambda-1|$.'' The corpus entry for P12 needs revision.


%=============================================================================
\section{SQMS Bound Reliability}
%=============================================================================

With $c_\psi = c$ confirmed:
\begin{itemize}
\item The psi-mode frequency $\Opsi = c k_{\rm cav} \sim 10^{10}$~rad/s.
  The SRF cavity is in the $\omega \sim \Opsi$ regime (not $\omega \gg
  \Opsi$), so the loss formula from W4i15 applies.
\item The radiation-$Q$ estimate $\Qpsi^{\rm local} \sim (kR)^3 \sim 31$
  is physically meaningful when $c_\psi = c$.
\item The SQMS bound scales as $\lbone < 1.56\ee{-9} \times
  \sqrt{\Qpsi/10^9}$. If $\Qpsi^{\rm local} = 31$, the bound
  tightens to $\lbone < 2.7\ee{-13}$.
\item The Q-ratio diagnostic ($\mathcal{R} = 0.49$) is ROBUST
  (confirmed by both agents).
\item \textbf{Open question}: Does the SRF cavity wall provide
  partial $\psi$-mode confinement (enhancing $\Qpsi^{\rm local}$
  above 31)? This depends on the boundary conditions for $\dpsiEM$
  at the Nb surface. If the psi-field is transparent to matter
  boundaries (which it should be, since $\psi$ couples gravitationally),
  then $\Qpsi^{\rm local} \sim 31$ and the bound tightens dramatically.
\end{itemize}

\textbf{Verdict}: SQMS bound reliable in structure. Absolute value
has an open $\sqrt{\Qpsi}$ factor. Range: $\lbone < [2.7\ee{-13},
1.56\ee{-9}]$. The Q-ratio shape measurement is the robust observable.


%=============================================================================
\section{Detection Programme}
%=============================================================================

The detection programme structure SURVIVES with $c_\psi = c$:
\begin{enumerate}
\item \textbf{Phase~0 pathfinder} (\$50K, 2 weeks): Go/no-go for
  Phase~I. UNAFFECTED (measures Q-ratio shape, independent of $\Qpsi$).
\item \textbf{Phase~I SQMS} (\$3.2M): Multi-frequency Q-ratio
  measurement. UNAFFECTED (shape measurement).
\item \textbf{Absolute bound}: If $\Qpsi^{\rm local} = 31$, the
  existing SQMS data already constrains $\lbone < 2.7\ee{-13}$,
  which is TIGHTER than the Phase~I target. This would make Phase~I
  a confirmation (or refutation) exercise, not a discovery.
\item \textbf{MEMS programme}: UNAFFECTED (uses passive observables,
  no $\Qpsi$ dependence).
\end{enumerate}


%=============================================================================
%=============================================================================
\part{Finding 5: Corpus Rewrite Errors}
\label{part:corpus}
%=============================================================================
%=============================================================================

The i16 corpus rewrite (MASTER\_FINDINGS\_CORPUS.md) was mandated to
integrate all findings through i15. Cross-checking against the i16
iteration tracker and the i16 agent outputs reveals 4 errors:

\begin{enumerate}
\item \textbf{P12 death cause}: Corpus says ``DEAD. Downgraded from
  THEORETICAL (W4i15). Multi-cavity superradiance KILLED by $c_\psi
  \sim 10^{-5}$~m/s.'' The i16 tracker flags ``P12: DEAD $\to$
  POTENTIALLY REVIVED (pending $c_\psi$ resolution)'' and Agent~4
  showed $c_\psi = c$ in the lab. P12 should be listed as DEAD with
  the corrected death cause (cooperativity too low at allowed
  $|\lambda-1|$), not the $c_\psi$ argument.

\item \textbf{Dust-branch $w$}: Corpus says ``$w \sim 10^{-37}$''
  in two places (Section~4, key quantities table). The i16 tracker
  explicitly states ``Dust-branch $w$: $10^{-37} \to 5\ee{-6}$''
  as a correction. The corpus was not updated with this i16 correction.

\item \textbf{$c_\psi$ value}: Corpus key quantities table says
  ``$c_\psi \sim 10^{-5}$~m/s, Dust branch (W4i14).'' The i16
  tracker states ``$c_\psi$: single value $\to$ regime-dependent
  (CONTRADICTORY between agents).'' At minimum, the corpus should
  flag the regime dependence. With the i17 resolution: ``$c_\psi = c$
  (laboratory, large-$\Delta$); $c_\psi \to 0$ (cosmological,
  small-$\Delta$).''

\item \textbf{``Definitively Closed'' list}: Entry ``Superluminal
  $\psi$-signaling (confirmed: $\psi$ propagates at $c_\psi \sim
  10^{-5}$~m/s, NOT $c$)'' is wrong. In the laboratory,
  $c_\psi = c$. The correct statement: ``Superluminal
  $\psi$-signaling impossible ($c_\psi \leq c$ in all regimes).''
  No-signaling is a stronger result when $c_\psi = c$ (information
  travels at most at $c$) than when $c_\psi \ll c$.
\end{enumerate}


%=============================================================================
%=============================================================================
\part{Summary: Programme Impact}
\label{part:summary}
%=============================================================================
%=============================================================================

\begin{tcolorbox}[colback=green!5!white, colframe=green!60!black,
  title=\textbf{i17 RESOLUTION SUMMARY}]

\begin{tabular}{@{}p{4cm}p{3cm}p{5cm}@{}}
\toprule
\textbf{Item} & \textbf{Pre-i17} & \textbf{Post-i17} \\
\midrule
$c_\psi$ (lab) & Contradictory & $c_\psi = c$ (confirmed) \\
$c_\psi$ (cosmo) & $10^{-5}$~m/s & $\to 0$ (dust branch; exact value background-dependent) \\
P1 & DEAD & DEAD (confirmed) \\
P8 & DEAD & DEAD (confirmed) \\
P12 & ``POTENTIALLY REVIVED'' & DEAD (cooperativity, not $c_\psi$) \\
SQMS bound & ``UNRELIABLE'' & Reliable but $\sqrt{\Qpsi}$ uncertain \\
Q-ratio & 0.49 (robust) & 0.49 (robust, confirmed) \\
$\Qpsi$ (local cavity) & Unknown & $\sim 31$ (radiation-limited at $c_\psi = c$; open question on wall confinement) \\
Corpus dust-branch $w$ & $10^{-37}$ (error) & $5\ee{-6}$ (needs correction) \\
\bottomrule
\end{tabular}

\textbf{Key resolution}: Agent~2 evaluated $K''$ at the background
$\Delta_{\rm bg}$, but the perturbation's own $\delta\Delta \gg
\Delta_{\rm bg}$ at SRF frequencies. The correct $c_\psi$ for
laboratory perturbations is $c$ (Agent~4). The ``contradiction'' was
a linearization scope error, not a physics disagreement.

\textbf{Action items for i18}:
\begin{enumerate}
\item Update corpus: P12 death cause, dust-branch $w$, $c_\psi$
  regime-dependence, superluminal entry.
\item Determine $\Qpsi^{\rm local}$ rigorously: does the psi-field
  see cavity walls? If not, $\Qpsi^{\rm local} = 31$ and the SQMS
  bound tightens $5700\times$.
\item Verify Agent~4's action normalization ($c^4/4$) discrepancy
  with the manuscript authors --- is it a presentation issue or a
  physics correction?
\end{enumerate}

\end{tcolorbox}


\end{document}
